% Part: first-order-logic 
% Chapter: axiomatic-deduction 
% Section: deduction-theorem

% pamriksan sipat-sipat provabilitas kang dibutuhake kanggo himpunan
% konsisten maksimal ing bab kelengkapan.

\documentclass[../../../include/open-logic-section]{subfiles}

\begin{document}

\iftag{FOL}
      {\olfileid{fol}{axd}{ded}}
      {\olfileid{pl}{axd}{ded}}

\olsection{Teorema Deduksi}

Kaya kang wis dideleng, nyuguhake !!{derivation}s ing sistem aksiomatik
iku rumit, lan !!{derivation}s kasebut bisa uga angel ditemokake.
Tinimbang nulis dhaptar !!{formula}s kang dawa kanthi rinci, lumrahe luwih
gampang mbuktekake yen !!{derivation}s kasebut ana kanthi nggunakake
sawetara asil prasaja. Kita wis netepake telung asil kaya mangkono:
\olref[ptn]{prop:reflexivity} nyebutake yen kita tansah bisa negesake
$\Gamma \Proves !A$ nalika ngerti yen $!A \in \Gamma$.
\olref[ptn]{prop:monotonicity} nyebutake yen, yen $\Gamma \Proves !A$,
mula uga $\Gamma \cup \{!B\} \Proves !A$. Lan
\olref[ptn]{prop:transitivity} menehi akibat yen, yen $\Gamma \Proves
!A$ lan $!A \Proves !B$, mula $\Gamma \Proves !B$. Ana asil
prasaja liyane, yaiku versi “meta” saka modus ponens:

\begin{prop}
  \ollabel{prop:mp} Yen $\Gamma \Proves !A$ lan $\Gamma \Proves !A
  \lif !B$, mula $\Gamma \Proves !B$.
\end{prop}

\begin{proof}
  Kita duwe $\{!A, !A \lif !B\} \Proves !B$:
  \begin{derivation}
    1. & $!A$ & Hip.\\
    2. & $!A \lif !B$ & Hip.\\
    3. & $!B$ & 1, 2, MP
  \end{derivation}
  Miturut \olref[ptn]{prop:transitivity}, $\Gamma \Proves !B$.
\end{proof}

Asil paling wigati kang bakal digunakake ing konteks iki yaiku
teorema deduksi:

\begin{thm}[Teorema Deduksi]
\ollabel{thm:deduction-thm} $\Gamma \cup \{!A\} \Proves !B$ yen lan
mung yen $\Gamma \Proves !A \lif !B$.
\end{thm}

\begin{proof}
Arah “yen” langsung lumaku. Yen $\Gamma \Proves !A \lif !B$, mula
uga $\Gamma \cup \{!A\} \Proves !A \lif !B$ miturut
\olref[ptn]{prop:monotonicity}. Uga, $\Gamma \cup \{!A\} \Proves !A$
miturut \olref[ptn]{prop:reflexivity}. Mula, miturut
\olref{prop:mp}, $\Gamma \cup \{!A\} \Proves !B$.

Kanggo arah “mung yen”, kita lumaku kanthi induksi miturut dawane
!!{derivation} kang formula pungkasane $!B$ saka $\Gamma \cup \{!A\}$.

Kanggo kasus dhasar induksi, kita mbuktekake pratelan kasebut kanggo
saben !!{derivation} kang dawane~$1$. !!^a{derivation} kang formula
pungkasane~$!B$ saka $\Gamma \cup \{!A\}$ lan dawane~$1$ mung dumadi
saka $!B$ dhewe; lan yen iku bener, $!B \in \Gamma \cup \{!A\}$ utawa
iku aksioma.
Cathetan koreksi sumber OLPL-027: sumber Inggris misahake B saka
tandha kaanggotaan, dene pamisahan kasus sabanjure mbutuhake
pratelan kaanggotaan kang pepak iki. Yen $!B \in \Gamma$ utawa iku
aksioma, mula $\Gamma \Proves !B$. Kita uga duwe $\Gamma
\Proves !B \lif (!A \lif !B)$ miturut \olref[prp]{ax:lif1}, lan
\olref{prop:mp} menehi $\Gamma \Proves !A \lif !B$. Yen $!B \in \{
!A\}$, mula $\Gamma \Proves !A \lif !B$, amarga !!{sentence}
pungkasan~$!A \lif !B$ padha karo $!A \lif !A$, lan kita wis
!!{derive} iki ing \olref[pro]{ex:identity}.

Kanggo langkah induktif, upamana !!a{derivation} kang formula pungkasane
$!B$ saka $\Gamma \cup \{!A\}$ lan langkah pungkasan~$!B$ kasebut
nduweni dhasar modus ponens. (Yen langkah iki minangka hipotesis utawa aksioma, mula
$!B \in \Gamma$, $!B \ident !A$, utawa $!B$ minangka aksioma, lan
pamikiran kang padha karo kasus dhasar induksi lumaku.) Cathetan sumber OLPL-736: ukara sumber iki mung nyakup hipotesis, aksioma, lan modus ponens. Nalika tag FOL aktif, aturan kuantor uga diidini; kasus kasebut dirembug ing perangan sabanjure, Teorema Deduksi mawa Kuantor. Banjur sawetara
langkah sadurunge ing !!{derivation} kasebut yaiku $!C \lif !B$ lan
$!C$, kanggo sawijining !!{formula}~$!C$, yaiku $\Gamma \cup \{!A\}
\Proves !C \lif !B$ lan $\Gamma \cup \{!A\} \Proves !C$, lan
!!{derivation}s kasebut luwih cekak, mula hipotesis induksi lumaku
kanggo kalorone. Dadi, kita nduweni:
\begin{align*}
  & \Gamma \Proves !A \lif (!C \lif !B); \\
  & \Gamma \Proves !A \lif !C.
\end{align*}
Nanging uga
\[
\Gamma \Proves (!A \lif (!C \lif !B)) \lif
((!A\lif !C)  \lif (!A \lif !B)),
\]
miturut \olref[prp]{ax:lif2}, lan rong aplikasi
\olref{prop:mp} menehi $\Gamma \Proves !A \lif !B$, kaya kang
dikarepake.
\end{proof}

Gatekna carane \olref[prp]{ax:lif1} lan \olref[prp]{ax:lif2} dipilih
mligi supaya Teorema Deduksi lumaku.

Ing ngisor iki sawetara kasunyatan migunani ngenani
!!{derivability}, kang ditinggalake minangka gladhen.

\begin{prop}
  \ollabel{prop:derivfacts}
\begin{enumerate}
\item $\Proves (!A \lif !B) \lif ((!B \lif !C)
  \lif (!A \lif !C))$; \ollabel{derivfacts:a}
  Cathetan koreksi sumber OLPL-028: siji kurung panutup kang ilang ing
  sumber Inggris dibalekake supaya formula iki imbang.
\item Yen $\Gamma \cup \{ \lnot !A\}
  \Proves \lnot !B$, mula $\Gamma \cup \{ !B\} \Proves
  !A$ (Kontraposisi); \ollabel{derivfacts:b}
\item  $\{ !A, \lnot!A\} \Proves
    !B$ (Ex Falso Quodlibet, Eksplosi); \ollabel{derivfacts:c}
\item  $\{ \lnot\lnot!A\} \Proves
  !A$ (Eliminasi Negasi Rangkep);\ollabel{derivfacts:d}
\item Yen $\Gamma \Proves \lnot\lnot!A$, mula $\Gamma \Proves
  !A$;\ollabel{derivfacts:e}
\end{enumerate}
\end{prop}

\tagprob{FOL}
\begin{prob}
  Buktekna \olref[fol][axd][ded]{prop:derivfacts}
\end{prob}
\tagendprob

\tagprob{notFOL}
\begin{prob}
  Buktekna \olref[pl][axd][ded]{prop:derivfacts}
\end{prob}
\tagendprob

\end{document}
